\section{Methods}
\label{sec:methods}

\subsection{Environment and task}
\label{sec:env}

We use XLand-MiniGrid~0.9.3 \citep{nikulin2024xland}, a JAX re-implementation of grid worlds
with objects and rules in the style of MiniGrid \citep{chevalierboisvert2023minigrid}, running on
CPU. Every episode takes place on a $7\times7$ grid built by the environment's procedural
generator from an integer \emph{map seed}; the seed fixes the walls and the positions of the
agent and the objects. The agent sees only an egocentric $5\times5$ crop in front of it: for
each cell, the tile type and its colour, scaled to $[0,1]$ and flattened to 50 inputs. It never
receives the goal, the rules or the global map. It has six discrete actions (move forward, turn
left, turn right, pick up, put down, toggle).

Task~B places two objects on the grid: a red ball, which is the goal, and a green key, which is a
distractor. The episode succeeds, and ends, when the agent holds the red ball; it is truncated
after 32 steps. The reward is the environment's default sparse success reward,
$1 - 0.9\,t/32$ at the successful step $t$ and $0$ otherwise. Because the view is partial, the
target is often out of sight, and a policy without memory can condition only on the current crop.
A random policy succeeds on 15.6\% of the validation maps and 14.1\% of the test maps.

\paragraph{Map panels.}
Maps are drawn from fixed, disjoint \emph{panels} of map seeds: 64 training maps (the only maps
seen during training, cycled episode by episode), 32 validation maps (learning curves and
checkpoint selection), 64 test maps (the reported score), plus a sealed final panel and a
small known panel that this paper does not use. Panels contain no duplicate layout within or
across panels: we hash each map's physical state (grid, agent position and direction, and the
contents of the agent's pocket) and remove repeats.

\subsection{Architectures}
\label{sec:arch}

Both agents share an encoder: a two-layer MLP with $\tanh$ activations, $50\to32\to32$. The
\textbf{MLP} agent reads its policy and value directly from the encoder and has no memory
(2,919 trainable parameters). The \textbf{LSTM} agent adds a single LSTM cell
\citep{hochreiter1997lstm} with 32 units on top of the encoder (11,367 parameters). Both use a
linear policy head over the six actions and a linear value head. The recurrent state is set to
zero at the start of every episode, in training and in evaluation. The two agents differ in size
by a factor of 3.9; \cref{sec:part2} controls for that difference.

\subsection{Training}
\label{sec:training}

Agents are trained with synchronous advantage actor-critic (A2C; \citealp{mnih2016a3c}) on a
single environment, with Adam \citep{kingma2015adam}, learning rate $10^{-3}$, rollouts of 16
steps, discount $\gamma=0.99$, value-loss coefficient $0.5$, entropy coefficient $0.01$ and
gradient-norm clipping at $0.5$. Gradients flow through the recurrent state within a rollout and
are cut at rollout boundaries. Following \citet{pardo2018timelimits}, the 32-step time limit is
treated as a truncation: the return bootstraps from the value of the last observation instead of
treating it as terminal. Each run lasts 200,000 environment interactions. These settings were
fixed by the study protocol before the pilot and are the same for both agents.

\subsection{Evaluation}
\label{sec:eval}

\paragraph{During training.}
Every 10,000 interactions (21 points per run, including step 0) the current weights are frozen,
saved as a checkpoint, and played once on each of the 32 validation maps. Actions are sampled from
the policy (not taken greedily); the sampling generator of each episode is seeded with its map
seed, so an evaluation is a deterministic function of the weights. This gives the learning curve.
At the end of training, the final weights are played once on each of the 64 test maps; this
\emph{final test success} is the primary score of a run.

\paragraph{Derived metrics.}
The \emph{validation AUC} is the trapezoidal area under the learning curve divided by the
training length, so it rewards learning early as well as learning well. The protocol's
\emph{acquisition rule} declares a task acquired when two consecutive validation points reach
at least $\max(0.75,\; s_0+0.2,\; s_{\text{rand}}+0.2)$, where $s_0$ is the success at step~0 and
$s_{\text{rand}}$ that of the random policy; the acquisition step is the step of the second such
point. A run therefore has a yes/no \emph{acquired} verdict.

\paragraph{Re-evaluation.}
To measure the randomness of evaluation itself, every saved checkpoint (all 21 curve points and
the final weights) is played again $R=3$ times with the weights fixed, using a new sampling seed
for each map and replica, $\text{map seed} + r\cdot 10{,}000{,}019$ for $r=1,2,3$ (replica~0 is
the original evaluation). We keep the outcome of every episode. Re-evaluation never changes the
run it reads; it verifies each checkpoint's SHA-256 digest before loading it.

\subsection{Seeds and preregistration}
\label{sec:seeds}

A run is identified by its architecture and its \emph{training seed}, which fixes the initial
weights, the order in which training maps are visited and the action-sampling stream during
training. Before any run started, a plan file fixed the study, the two architectures, the task,
the budget, $R$, and the seed rule: the three pilot seeds (7, 19, 31) followed by 27 seeds drawn
from NumPy's \texttt{SeedSequence(20261018)} modulo $10^6$, skipping repeats. Both architectures
are trained on the same 30 seeds, so every comparison is paired by seed. The plan, and a hash of
it, are recorded in the chain's journal; a changed plan cannot continue an existing chain.

\subsection{Separating seed, map and evaluation noise}
\label{sec:gtheory}

For one architecture and one metric, let $y_{smr}$ be the outcome of seed $s$ on map $m$ in
re-evaluation replica $r$ (for final test success, $y_{smr}\in\{0,1\}$; for the AUC, the area
under that map's curve in that replica). We use the crossed random-effects model of
generalizability theory \citep{cronbach1972dependability,brennan2001generalizability}
\begin{equation}
  y_{smr} = \mu + a_s + b_m + (ab)_{sm} + e_{smr},
  \label{eq:model}
\end{equation}
with independent zero-mean effects of variance $\sigma^2_{\text{seed}}$ (real differences between
training seeds), $\sigma^2_{\text{map}}$ (some maps are harder for everyone),
$\sigma^2_{\text{seed}\times\text{map}}$ (a seed finds some maps easier than other seeds do) and
$\sigma^2_{\text{noise}}$ (the same weights on the same map give a different outcome under another
sampling seed). With $S$ seeds, $M$ maps and $R$ replicas, the balanced ANOVA estimators
\citep{searle1992variance} are
\begin{align*}
  \hat\sigma^2_{\text{noise}} &= \mathrm{MS}_E, &
  \hat\sigma^2_{\text{seed}\times\text{map}} &= (\mathrm{MS}_{SM}-\mathrm{MS}_E)/R,\\
  \hat\sigma^2_{\text{seed}} &= (\mathrm{MS}_S-\mathrm{MS}_{SM})/(MR), &
  \hat\sigma^2_{\text{map}} &= (\mathrm{MS}_M-\mathrm{MS}_{SM})/(SR),
\end{align*}
where $\mathrm{MS}$ are the usual mean squares; negative estimates are reported as zero. The
acquisition verdict has one value per seed and replica, so for it we use the one-way version
($M=1$, map terms absent). Paired differences between architectures (LSTM minus MLP on the same
seed, map and sampling seed) are analysed with the same model. Intervals for the components are
percentile intervals from 1,000 bootstrap resamples of seeds \citep{efron1993bootstrap}.

\paragraph{From components to designs (D study).}
The variance of an architecture's mean score over $n_s$ seeds, $n_m$ maps and $n_r$ evaluations is
\begin{equation}
  \operatorname{Var}(\bar y) = \frac{\sigma^2_{\text{seed}}}{n_s}
    + \frac{\sigma^2_{\text{seed}\times\text{map}}}{n_s n_m}
    + \frac{\sigma^2_{\text{noise}}}{n_s n_m n_r}
    \;\Big[+\frac{\sigma^2_{\text{map}}}{n_m}\Big],
  \label{eq:dstudy}
\end{equation}
where the bracketed term applies only to claims about new maps; on a fixed panel shared by all
runs it is common to every arm and cancels in a comparison. From \cref{eq:dstudy} we report
(i)~the smallest $n_s$ whose 95\% interval half-width $1.96\sqrt{\operatorname{Var}(\bar y)}$ is
below a target, and (ii)~for a paired difference with per-seed variance
$V=\sigma^2_{\text{seed}}+\sigma^2_{\text{seed}\times\text{map}}/n_m+\sigma^2_{\text{noise}}/(n_m n_r)$,
the seeds a two-sided test at $\alpha=0.05$ needs to find a true gap $\delta$ with power 0.8,
$n_s=\lceil(z_{0.975}+z_{0.8})^2 V/\delta^2\rceil$ (normal approximation, minimum 2). Finally,
the \emph{noise share} of a metric is the fraction of the between-seed variance of a run's score
from a single evaluation ($n_m$ maps, $n_r=1$) that is due to evaluation noise,
$(\sigma^2_{\text{noise}}/n_m)\,/\,V$. All computations are in
\texttt{ucl study-variance}, which we checked on synthetic data with known components.

\subsection{Execution}
\label{sec:execution}

Runs execute one at a time as a chain of events: the end of a run starts its re-evaluation, and
the end of a re-evaluation starts the next run. Every event is appended to the chain's journal; a
failed or interrupted step stops the chain and is shown as such until it is repeated by hand with
the same seed. Two runs were stopped and repeated with the same seed, their partial outputs kept:
one by a planned restart of the machine, and one by the runner's resource guard, which halts a
run when available memory falls below 400\,MiB. Everything ran on
one laptop CPU (Intel Core i7-12700H, WSL2, PyTorch~2.11 with two threads, JAX~0.4.38 on CPU).
In Part~I a run took 6.4 minutes for the LSTM and 4.7 for the MLP, and its re-evaluation about
one minute; the 60 runs used 12M training interactions and about 6.5 hours. Part~II's 128 runs
(including the learning-rate search and the mechanism arms) used 38.4M interactions and about
24 hours.
